\textbf{Phần 1. Trắc nghiệm lựa chọn}
\setcounterpageref{ex}{0}

\begin{ex} %[5H8N45]
Thể tích của một hình là
\choice
{Diện tích toàn bộ bề mặt bên ngoài của hình đó}
{\True Độ lớn của khoảng không gian mà hình đó chiếm chỗ}
{Tổng độ dài tất cả các cạnh của hình đó}
{Khối lượng của hình đó khi đem cân}
\loigiai{
Thể tích của một hình là độ lớn của khoảng không gian mà hình đó chiếm chỗ.\\
\textbf{Đáp án của câu hỏi là B}
}
\end{ex}

\begin{ex} %[5H8N45]
Cho hình hộp chữ nhật A nằm hoàn toàn bên trong hình hộp chữ nhật B. Kết luận nào sau đây là đúng?
\choice
{Thể tích hình A lớn hơn thể tích hình B}
{Thể tích hình A bằng thể tích hình B}
{Không so sánh được thể tích của hai hình}
{\True Thể tích hình A bé hơn thể tích hình B}
\loigiai{
Vì hình A nằm hoàn toàn trong hình B nên thể tích hình A bé hơn thể tích hình B.\\
\textbf{Đáp án của câu hỏi là D}
}
\end{ex}

\begin{ex} %[5H8N45]
Một hình khối $M$ được ghép bởi $12$ khối lập phương nhỏ bằng nhau. Hình khối $N$ được ghép bởi $15$ khối lập phương nhỏ như thế. Kết luận nào đúng?
\choice
{\True Thể tích hình $N$ lớn hơn thể tích hình $M$}
{Thể tích hình $M$ lớn hơn thể tích hình $N$}
{Thể tích hai hình bằng nhau}
{Thể tích hình $M$ gấp đôi thể tích hình $N$}
\loigiai{
Hình $M$ gồm $12$ khối lập phương nhỏ, hình $N$ gồm $15$ khối lập phương nhỏ nên thể tích hình $N$ lớn hơn thể tích hình $M$.\\
\textbf{Đáp án của câu hỏi là A}
}
\end{ex}

\begin{ex} %[5H8H45]
Hình C gồm $8$ khối lập phương cạnh $1\text{ cm}$. Hình D gồm $8$ khối lập phương cạnh $1\text{ cm}$ xếp theo hình dạng khác. So sánh thể tích của hình C và hình D:
\choice
{Hình C có thể tích lớn hơn hình D}
{Hình D có thể tích lớn hơn hình C}
{Không thể so sánh được}
{\True Hình C và hình D có thể tích bằng nhau}
\loigiai{
Hai hình được ghép từ cùng số lượng các khối lập phương như nhau thì có thể tích bằng nhau.\\
\textbf{Đáp án của câu hỏi là D}
}
\end{ex}

\begin{ex} %[5H8H45]
Hình $P$ được ghép từ hai hình khối $Q$ (gồm $6$ khối lập phương đơn vị) và $R$ (gồm $9$ khối lập phương đơn vị). Thể tích của hình $P$ bằng
\choice
{\True $15$ khối lập phương đơn vị}
{$3$ khối lập phương đơn vị}
{$54$ khối lập phương đơn vị}
{$18$ khối lập phương đơn vị}
\loigiai{
Thể tích hình $P$ bằng tổng thể tích của hai hình $Q$ và $R$: $6 + 9 = 15$ khối lập phương đơn vị.\\
\textbf{Đáp án của câu hỏi là A}
}
\end{ex}

\begin{ex} %[5H8H45]
Một hình lập phương lớn được ghép từ $27$ khối lập phương nhỏ cạnh $1\text{ cm}$. Người ta lấy ra $3$ khối lập phương nhỏ ở $3$ góc. Thể tích của hình còn lại là
\choice
{$27\text{ khối}$}
{$21\text{ khối}$}
{\True $24\text{ khối}$}
{$18\text{ khối}$}
\loigiai{
Thể tích còn lại: $27 - 3 = 24$ khối lập phương nhỏ.\\
\textbf{Đáp án của câu hỏi là C}
}
\end{ex}

\begin{ex} %[5H8V45]
Người ta dùng các khối lập phương cạnh $1\text{ cm}$ để xếp thành một hình hộp chữ nhật có chiều dài $5\text{ cm}$, chiều rộng $4\text{ cm}$ và chiều cao $3\text{ cm}$. Số khối lập phương nhỏ cần dùng là
\choice
{$12$ khối}
{\True $60$ khối}
{$20$ khối}
{$15$ khối}
\loigiai{
Mỗi lớp có $5 \times 4 = 20$ khối. Có $3$ lớp nên tổng số khối là $20 \times 3 = 60$ khối.\\
\textbf{Đáp án của câu hỏi là B}
}
\end{ex}

\begin{ex} %[5H8V45]
Bạn Nam có $48$ khối lập phương nhỏ bằng nhau. Nam muốn xếp tất cả các khối đó thành một hình hộp chữ nhật. Cách xếp nào sau đây KHÔNG thể thực hiện được?
\choice
{Dài $4$, rộng $3$, cao $4$}
{Dài $6$, rộng $4$, cao $2$}
{\True Dài $5$, rộng $3$, cao $3$}
{Dài $8$, rộng $3$, cao $2$}
\loigiai{
Vì $5 \times 3 \times 3 = 45 \neq 48$ nên không thể xếp thành hình hộp chữ nhật có kích thước $5, 3, 3$.\\
\textbf{Đáp án của câu hỏi là C}
}
\end{ex}

\begin{ex} %[5H8V45]
Một hình lập phương lớn có cạnh $3\text{ dm}$ được ghép bởi các hình lập phương nhỏ cạnh $1\text{ dm}$. Người ta sơn $6$ mặt ngoài của hình lập phương lớn. Số hình lập phương nhỏ không được sơn mặt nào là
\choice
{$0$ hình}
{$6$ hình}
{$8$ hình}
{\True $1$ hình}
\loigiai{
Số hình lập phương nhỏ ở tâm không sơn mặt nào là $(3 - 2) \times (3 - 2) \times (3 - 2) = 1$ hình.\\
\textbf{Đáp án của câu hỏi là D}
}
\end{ex}

\begin{ex} %[5H8V45]
Thả một hòn đá chìm hoàn toàn vào một cốc nước hình trụ đang chứa nước. Mực nước trong cốc dâng lên từ vạch $120\text{ ml}$ đến vạch $165\text{ ml}$. Thể tích của hòn đá là
\choice
{$285\text{ ml}$}
{$165\text{ ml}$}
{\True $45\text{ ml}$}
{$120\text{ ml}$}
\loigiai{
Thể tích hòn đá bằng thể tích nước dâng lên: $165 - 120 = 45\text{ (ml)}$.\\
\textbf{Đáp án của câu hỏi là C}
}
\end{ex}

\begin{ex} %[5H8V45]
Có $64$ khối lập phương nhỏ cạnh $1\text{ cm}$ được xếp lại thành một hình lập phương lớn. Cạnh của hình lập phương lớn là
\choice
{$2\text{ cm}$}
{\True $4\text{ cm}$}
{$6\text{ cm}$}
{$8\text{ cm}$}
\loigiai{
Vì $4 \times 4 \times 4 = 64$ nên cạnh hình lập phương lớn là $4\text{ cm}$.\\
\textbf{Đáp án của câu hỏi là B}
}
\end{ex}

\begin{ex} %[5H8C45]
Một khối gỗ hình hộp chữ nhật được ghép từ $120$ khối lập phương nhỏ cạnh $1\text{ cm}$. Người ta sơn đỏ toàn bộ mặt ngoài của khối gỗ rồi tháo rời các khối nhỏ ra. Biết khối gỗ có kích thước $4\text{ cm} \times 5\text{ cm} \times 6\text{ cm}$. Số khối lập phương nhỏ được sơn đúng $2$ mặt là
\choice
{\True $36$ khối}
{$24$ khối}
{$48$ khối}
{$16$ khối}
\loigiai{
Số khối sơn 2 mặt nằm trên các cạnh (trừ 8 đỉnh):\\
$4 \times (4 - 2) + 4 \times (5 - 2) + 4 \times (6 - 2) = 4 \times 2 + 4 \times 3 + 4 \times 4 = 8 + 12 + 16 = 36$ (khối).\\
\textbf{Đáp án của câu hỏi là A}
}
\end{ex}

\textbf{Phần 2. Trắc nghiệm đúng sai}
\setcounterpageref{ex}{0}

\begin{ex} %[5H8H45]
Cho hình hộp chữ nhật $A$ gồm $24$ khối lập phương đơn vị và hình hộp chữ nhật $B$ gồm $30$ khối lập phương đơn vị. Xét tính đúng sai:
\choiceTF
{\True Thể tích hình B lớn hơn thể tích hình A.}
{Thể tích hình A bằng $\dfrac{5}{4}$ thể tích hình B.}
{\True Thể tích hình A bằng $\dfrac{4}{5}$ thể tích hình B.}
{\True Khi ghép hình A và hình B lại ta được một hình mới có thể tích bằng $54$ khối lập phương đơn vị.}
\loigiai{
\begin{itemize}
    \item $30 > 24$ nên thể tích B lớn hơn thể tích A. Đúng.
    \item Tỉ số A so với B là $24 : 30 = \dfrac{4}{5} \neq \dfrac{5}{4}$. Ý 2 Sai, Ý 3 Đúng.
    \item Ghép lại: $24 + 30 = 54$ khối. Đúng.
\end{itemize}
\begin{itemize}
    \item a) \textbf{Đúng}.
    \item b) \textbf{Sai}.
    \item c) \textbf{Đúng}.
    \item d) \textbf{Đúng}.
\end{itemize}
}
\end{ex}

\begin{ex} %[5H8V45]
Xếp các khối lập phương cạnh $1\text{ cm}$ thành một hình lập phương lớn có cạnh $4\text{ cm}$, sau đó sơn tất cả các mặt ngoài. Xét tính đúng sai:
\choiceTF
{\True Cần dùng tất cả $64$ khối lập phương nhỏ.}
{\True Có đúng $8$ khối lập phương nhỏ được sơn $3$ mặt.}
{Có $16$ khối lập phương nhỏ không được sơn mặt nào.}
{\True Có $24$ khối lập phương nhỏ được sơn đúng $1$ mặt.}
\loigiai{
\begin{itemize}
    \item Số khối: $4 \times 4 \times 4 = 64$ khối. Đúng.
    \item Số khối sơn 3 mặt ở 8 góc $= 8$ khối. Đúng.
    \item Số khối không sơn: $(4-2) \times (4-2) \times (4-2) = 8$ khối $\neq 16$ khối. Sai.
    \item Số khối sơn 1 mặt ở 6 mặt: $6 \times (4-2) \times (4-2) = 6 \times 4 = 24$ khối. Đúng.
\end{itemize}
\begin{itemize}
    \item a) \textbf{Đúng}.
    \item b) \textbf{Đúng}.
    \item c) \textbf{Sai}.
    \item d) \textbf{Đúng}.
\end{itemize}
}
\end{ex}

\textbf{Phần 3. Trả lời ngắn}
\setcounterpageref{ex}{0}

\begin{ex} %[5H8H45]
Hình khối $E$ được ghép bởi $18$ khối lập phương nhỏ cạnh $1\text{ cm}$, hình khối $F$ được ghép bởi $25$ khối lập phương nhỏ như thế. Hỏi thể tích hình $F$ lớn hơn thể tích hình $E$ bao nhiêu xăng-ti-mét khối?
\shortans{7}
\loigiai{
\textbf{Đáp số:} $7$\\
$25 - 18 = 7\text{ cm}^3$.
}
\end{ex}

\begin{ex} %[5H8H45]
Người ta xếp các khối lập phương nhỏ cạnh $1\text{ cm}$ thành hình hộp chữ nhật có kích thước dài $6\text{ cm}$, rộng $4\text{ cm}$, cao $5\text{ cm}$. Hỏi cần dùng tất cả bao nhiêu khối lập phương nhỏ?
\shortans{120}
\loigiai{
\textbf{Đáp số:} $120$\\
$6 \times 4 \times 5 = 120$ khối.
}
\end{ex}

\begin{ex} %[5H8V45]
Một hình lập phương lớn cạnh $3\text{ cm}$ được ghép từ các khối lập phương nhỏ cạnh $1\text{ cm}$. Sơn màu đỏ toàn bộ bên ngoài hình lập phương lớn. Hỏi có bao nhiêu khối lập phương nhỏ được sơn đúng $1$ mặt?
\shortans{6}
\loigiai{
\textbf{Đáp số:} $6$\\
Mỗi mặt có $1$ khối ở giữa: $6 \times 1 = 6$ khối.
}
\end{ex}

\begin{ex} %[5H8V45]
Thả chìm một quả cầu sắt vào bình chia độ chứa $250\text{ ml}$ nước thì nước dâng lên mức $380\text{ ml}$. Tính thể tích quả cầu theo mi-li-lít.
\shortans{130}
\loigiai{
\textbf{Đáp số:} $130$\\
$380 - 250 = 130\text{ ml}$.
}
\end{ex}

\begin{ex} %[5H8V45]
Để xếp một hình lập phương có cạnh $5\text{ cm}$ bằng các khối lập phương nhỏ cạnh $1\text{ cm}$, ta cần dùng bao nhiêu khối lập phương nhỏ?
\shortans{125}
\loigiai{
\textbf{Đáp số:} $125$\\
$5 \times 5 \times 5 = 125$ khối.
}
\end{ex}

\begin{ex} %[5H8C45]
Người ta xếp $216$ khối lập phương nhỏ cạnh $1\text{ cm}$ thành một hình lập phương lớn, sau đó sơn tất cả các mặt bên ngoài. Hỏi có bao nhiêu khối lập phương nhỏ không được sơn mặt nào?
\shortans{64}
\loigiai{
\textbf{Đáp số:} $64$\\
Cạnh hình lập phương lớn là $6\text{ cm}$ vì $6 \times 6 \times 6 = 216$.\\
Số khối không sơn mặt nào: $(6 - 2) \times (6 - 2) \times (6 - 2) = 4 \times 4 \times 4 = 64$ khối.
}
\end{ex}
